% !TEX root = ../main.tex

\chapter{Introduction}

\section{Motivation: Efficiency Bottlenecks in BSP}

The Bulk Synchronous Parallel (BSP) paradigm powers some of our most critical computational tools,
from large-scale scientific simulations in astrophysics and climate science to distributed training
of foundation models in machine learning.
These workloads---consisting of large parallel computations interspersed with I/O and global
synchronization---currently operate at the scale of tens of thousands of accelerator
nodes~\cite{Das2023:LargeScale,Stock2024:LargeScale,Gratt2024:Llama3Scale,:LlamaScale} and continue
to grow with time~\cite{Rahma2024:MLScale}.
As these applications scale, their efficiency has datacenter-level impacts on both cost and
time-to-discovery.

The tightly synchronized nature of BSP creates two unique challenges that strain modern computing
infrastructure.
On the data path, bursts of checkpointed state create massive I/O demands on the storage and
analytics infrastructure~\cite{Eisen2022:Checkpoint}.
On the compute path, frequent global synchronization amplifies the system's sensitivity to
execution variability~\cite{Hoefl2010:Noise,Petri2003:Missing,Yildi2016:IONoise}.
Performance bottlenecks that would be manageable in loosely coupled systems, such as minor OS-level
interruptions or thermal throttling, become critical scaling barriers when every worker must wait
for the slowest participant.

We now describe the status quo and outstanding efficiency challenges along the data and compute
paths.

\paragraph{Data Management}
The explosive growth in scientific data volumes has made in-situ analytics an essential strategy
for reducing I/O overhead and accelerating time-to-insight.
While widely used for visualization and occasionally adopted for data transformation or staging,
current in-situ systems are fragmented, limited, and purpose-built.
Most frameworks hardcode specific operations or assume fixed data layouts, with no support for
composing multiple indexing strategies at ingest time.
Existing systems like FastQuery build auxiliary indices such as bitmaps, while others like DeltaFS
perform high-throughput hash partitioning suitable for point queries—but neither supports efficient
range queries.
No general-purpose mechanism exists for building query-efficient layouts dynamically and
incrementally during ingestion.
As a result, range query performance still hinges on costly post-processing sorts, undermining the
core premise of in-situ data processing.

\paragraph{Compute Efficiency}
The parallel inefficiency of BSP codes is widely ascribed to load imbalance, but there has been
little systematic effort to characterize its sources or validate mitigation strategies.
The standard approach—offline trace analysis—is both cumbersome and disconnected from real-time
system context.
Traces are large, opaque, and typically analyzed long after execution, often by researchers without
access to the hardware environment that produced them.
Profiling tools, while lightweight, expose only coarse-grained summaries and miss transient effects
that may drive synchronization delays.
As a result, performance tuning remains ad-hoc and reactive, lacking the infrastructure for
targeted, context-aware introspection during execution.

\section{Our Work: A Comprehensive Study of BSP Efficiency}

Our work developed systems and algorithms to advance workflow efficiency of parallel codes on
both the compute and data paths.
These systems demonstrated that significant efficiency gains were possible and revealed key
insights for making these wins accessible on production clusters.
This experience led us to develop ORCA, a general-purpose approach that enables efficiency wins
in production environments.
We describe each system in turn, highlighting the contributions and insights that collectively
motivated our unified approach.

\subsection{Data Path: Streaming Data Reorganization}

\paragraph{Motivation}
Our work was motivated by the need for write-efficient ingestion techniques that for scientific
data analyses.
Analyses on continuous physical properties such as wavefronts or energy bands require efficient
range queries.
All existing solutions for efficient range queries require reordering of on-disk data via
postprocessing.
Not only would an in-situ ingestion system deliver efficiency gains, it would also complete the
toolkit of fundamental in-situ data organization primitives.

\paragraph{Challenges}
Our goal was to develop a purely streaming system that would never reorder data once written to
disk, without making any assumptions about data distribution characteristics.
Our characterization study found scientific data distributions to be highly skewed and evolving
over time.
Any static range partitioning scheme would become imbalanced, creating write bottlenecks and query
hotspots.
Our system needed to ingest data with unknown, dynamic, a nd evolving data distributions, maintain
full utilization of the storage system, and scale to thousands of cores with minimal application
runtime overhead.

\paragraph{Our Contribution}
CARP (Continuous Adaptive Range Partitioner), an adaptive partitioner that adapts to arbitrary key
distributions and partitions them in a single streaming pass.
CARP provides query latencies comparable to those of a full sort without any ingestion overhead,
making it 5$\times$ faster than prior work.

% \paragraph{Lessons Learned} Analytics capabilities are fragmented across specialized tools,
% requiring database-style modularity that decouples analysis from data layout. Adaptive control
% loops are essential for dynamic workloads, and in-situ processing enables "good enough" solutions
% that outperform "perfect" offline approaches.

\subsection{Compute Path: Understanding and Mitigating Load Imbalance}

\paragraph{Motivation}
Poor utilization in adaptive mesh codes a well-known problem.
We wanted to create load-balancing mechanisms that would incorporate fine-grained runtime telemetry
from these codes, and incorporate them into a placement policy that conferred meaningful runtime
gains.

\paragraph{Challenges}
Our efforts to collect and interpret telemetry, and create models of code performance that could be
exploited in a placement policy encountered several persistent challenges.

\begin{enumerate}
  \item Raw telemetry measurements did not correlate with expected behavior, forcing deep investigation of
        cross-stack performance anomalies, and revealing issues with hardware, network stack parameters, operation
        ordering etc.
        These constituted runtime noise that had to be eliminated before telemetry and plpacement effects
        could be reliably observed.
  \item Identifying, isolating, and tracing transient performance issues in a large parallel code required
        collecting and analyzing a large volume of telemetry.
        Existing performance tools were unhelpful: profiling tools did not provide the required
        granuliarty, while traces were slow and not analyzable.
  \item Need efficient placement computation on NP-hard problem within tight timing constraints.
\end{enumerate}

\paragraph{Our Contribution}
A placement policy (CPLX) with tunable control over two placement tradeoffs: compute load balance
and communication locality.
By balancing these two objectives, CPLX provides upto 21.6\% runtime improvement over optimized
baselines.

% \paragraph{Lessons Learned} \begin{enumerate} \item "Load imbalance" is a misnomer - the real
% problem is cross-stack variability masquerading as algorithmic issues. \item Context-specific
% solutions are required as exported traces lose essential insight. \item Programmable telemetry is
% essential, and ad-hoc analytics pipelines proved more effective than standard tools. \item
% Real-time feedback is needed for effective diagnosis and tuning. \end{enumerate}

\subsection{ORCA: Transferring Lessons Learned to General-Purpose Observability}

\paragraph{Motivation}
While our previous efforts were successful in demonstrating the potential for efficiency gains on
both data and compute paths, realizing these gains on production clusters remained elusive.
Realizing data path efficiency requires a general-purpose data management system that can use the
best combination of indexes for a given analytics workload.
Realizing compute path efficieny is even more challenging: our work was conducted on an in-house
research cluster with no operational, security, or vendor constraints.
Realizing these benefits in production clusters requires efficient workflows to trace variability
to its root cause, which may or may not be amenable to algorithmic solutions.

\paragraph{Our Approach}
We are developing a general-purpose dataflow framework for cluster observability using SQL
semantics and Arrow.
ORCA (Observability with Realtime Control and Aggregation) uses a tree-based overlay network with
distributed query planning and a async control path.

\paragraph{Challenges}
No existing template exists for cluster-scale streaming analytics.
The system requires complex distributed coordination with multiple interacting components while
meeting strict performance requirements to avoid impacting application critical paths.

\section{Thesis Statement}

\begin{quote}
  This thesis asserts that overcoming the primary data and compute bottlenecks in large-scale BSP
  applications requires a paradigm shift from specialized, siloed tools to a unified, online
  \emph{sense-making} engine, and demonstrates that the architectural foundation for this shift is a
  framework that treats both in-situ analytics and performance diagnostics as a single problem of
  composing query-driven dataflows.
\end{quote}

\subsection{Potential Impact}

\paragraph{Immediate Impact}
Application-level observability becomes a first-class capability.
ORCA replaces brittle traces and low-resolution profilers with a lightweight, always-on,
query-driven system that supports real-time diagnosis and runtime tuning.
Developers can rapidly inspect distributed behavior, toggle probes, and react to anomalies without
modifying application code or halting execution.

\paragraph{Longer-Term Impact}
ORCA evolves into a flexible substrate for in-situ analytics and adaptive control.
Its open architecture supports programmable analysis pipelines, integration of diverse telemetry
sources, and runtime composition of custom indexing and aggregation strategies.
Over time, it can subsume visualization, introspection, and control workflows into a unified
framework that prioritizes insight per unit of data moved, enabling principled trade-offs between
data fidelity, performance, and resource usage.

\section{Document Structure}

The remainder of this document is organized as follows.
